\documentclass[11pt,a4paper]{article}

\usepackage[margin=2.7cm]{geometry}
\usepackage{amsmath}
\usepackage{amssymb}
\usepackage{booktabs}
\usepackage{hyperref}

\title{The Briggs model in \texttt{deamSim}:\\ what was wrong and what the fix does}
\author{}
\date{}

\newcommand{\Ol}{O_{5'}}

\begin{document}
\maketitle

\noindent
This note works through the two corrections made to the \texttt{-damage} option of
\texttt{deamSim}. It is written so that every step can be checked by hand; the only
probability facts used are the geometric distribution and independence.

\section{Setting up}

A fragment of length $L$. We follow one strand of the duplex, the one that will be
sequenced, and number its bases $n = 1, 2, \dots, L$ starting from its own $5'$ end.

Four parameters, as in Briggs et al.\ (2007):
\begin{center}
\begin{tabular}{ll}
\toprule
$\lambda$   & geometric parameter of the overhang length \\
$\delta_s$  & probability that a C in a single-stranded region is deaminated \\
$\delta_d$  & probability that a C in a double-stranded region is deaminated \\
$\nu$       & nick frequency per base \\
\bottomrule
\end{tabular}
\end{center}

\noindent
Their maximum likelihood estimates on the Neandertal data (their Table~1) are
$\nu = 0.024$, $\lambda = 0.36$, $\delta_d = 0.0097$, $\delta_s = 0.68$, and we will
use these throughout as a worked example.

Throughout, ``geometric with parameter $\lambda$'' means the version supported on
$\{0, 1, 2, \dots\}$,
\begin{equation}
\Pr(G = k) \;=\; \lambda (1-\lambda)^k , \qquad k = 0, 1, 2, \dots
\label{eq:geom}
\end{equation}
This is what \texttt{std::geometric\_distribution} gives, and what NGSNGS's
\texttt{Random\_geometric\_k} gives, so the two codebases agree on what $\lambda$
means. Its mean is $(1-\lambda)/\lambda$, which at $\lambda = 0.36$ is $1.78$
nucleotides --- matching the ``1.6--1.8 nucleotides'' quoted in the paper. Good; we
have the right parameterization.

The one quantity we will keep reaching for is the \emph{survival function}
\begin{equation}
\Pr(G \geq n) \;=\; \sum_{k \geq n} \lambda(1-\lambda)^k \;=\; (1-\lambda)^n .
\label{eq:surv}
\end{equation}

\section{Correction 1: the factor of one half}

\subsection{Where the one half comes from}

Consider the left end of the duplex. Both strands broke, at nearby but not identical
places. Two things can have happened:

\begin{itemize}
\item The strand we are going to sequence sticks out. Its $5'$ end is unpaired: a
      \textbf{$5'$ overhang}, single-stranded, free to deaminate.
\item The \emph{other} strand sticks out. Then our strand is the recessed one, and the
      protruding piece is the other strand's $3'$ end: a \textbf{$3'$ overhang}.
\end{itemize}

There is no reason for one to be more common than the other --- Briggs et al.\ say
``we assume overhangs are equally likely to extend either the $5'$ or $3'$ strand''
--- so each happens with probability $\tfrac12$.

Now library preparation. T4 polymerase blunt-end repair does two different things to
these two cases:
\begin{itemize}
\item a $5'$ overhang is \emph{filled in}: the strand itself survives intact, uracils
      and all, and those uracils are later read as T;
\item a $3'$ overhang is \emph{chewed back} by the $3' \!\to\! 5'$ exonuclease: the
      protruding piece is destroyed, and any damage it carried is gone forever.
\end{itemize}

So only half of the ends leave a single-stranded region we can ever observe. This is
exactly what the paper says in the main text: the end repair step ``eliminates all
$3'$-overhanging ends and preserves only $5'$-overhanging ends''.

\subsection{The resulting distribution}

Let $\Ol$ be the length of the \emph{observable} single-stranded region at the $5'$
end of our strand. Writing $B \sim \mathrm{Bernoulli}(\tfrac12)$ for ``this end was a
$5'$ overhang'' and $G$ for the geometric of~\eqref{eq:geom},
\begin{equation}
\Ol \;=\; B \cdot G .
\end{equation}
Explicitly,
\begin{equation}
\Pr(\Ol = 0) \;=\; \tfrac12 + \tfrac12 \lambda ,
\qquad
\Pr(\Ol = k) \;=\; \tfrac12\, \lambda (1-\lambda)^k \quad (k \geq 1).
\label{eq:mixture}
\end{equation}
Half a point mass at zero, half a geometric. And the quantity that will do all the
work, using~\eqref{eq:surv}:
\begin{equation}
\boxed{\;\Pr(\Ol \geq n) \;=\; \tfrac12 (1-\lambda)^n \;}
\qquad n \geq 1 .
\label{eq:survhalf}
\end{equation}

\medskip
\noindent
\textbf{What \texttt{deamSim} did instead.} It drew a geometric at both ends, with no
Bernoulli in front:
\begin{verbatim}
    int overhang5p = overhang(generator);
    int overhang3p = overhang(generator);
\end{verbatim}
that is, $\Ol = G$ and therefore $\Pr(\Ol \geq n) = (1-\lambda)^n$. Twice
\eqref{eq:survhalf}, at every single $n$.

\subsection{What that does to the damage}

Position $n$ is single-stranded exactly when the overhang reaches it, i.e.\ when
$\Ol \geq n$. Conditioning on that:
\begin{equation}
\Pr(\text{C}\!\to\!\text{T at } n)
 \;=\; \underbrace{\Pr(\Ol \geq n)\, \delta_s}_{\text{single-stranded}}
 \;+\; \underbrace{\Pr(\Ol < n)\, \Pr(\text{$n$ upstream of the nick})\, \delta_d}_{\text{double-stranded}} .
\label{eq:perpos}
\end{equation}
Since $\delta_d \approx 0.01$ while $\delta_s \approx 0.68$, near the end of the
molecule the first term is what matters. Substituting the two survival functions into
the first term of~\eqref{eq:perpos}:
\begin{align}
\text{Briggs:}    \qquad & \tfrac12 (1-\lambda)^n\, \delta_s , \\
\text{\texttt{deamSim} (old):} \qquad & (1-\lambda)^n\, \delta_s .
\end{align}
A clean factor of two at every position. At $n = 1$ with the paper's estimates:
\begin{equation}
\tfrac12 (1 - 0.36)(0.68) = 0.218
\qquad \text{versus} \qquad
(1 - 0.36)(0.68) = 0.435 .
\end{equation}

The same argument runs at the other end. There, the observable single-stranded region
belongs to the \emph{complementary} strand, whose deaminated Cs get copied into our
strand as A during the fill-in, which is why that end shows G$\to$A rather than
C$\to$T. But it is the same geometry and the same $\tfrac12$.

\subsection{Why this is not just a badly chosen $\lambda$}

The natural hope is that the old code was simply using a different $\lambda$. It was
not, and here is the one-line proof.

Suppose some $\lambda'$ made the old (pure geometric) model agree with the correct one
at every position. Then from~\eqref{eq:survhalf} we would need
\begin{equation}
(1-\lambda')^n \;=\; \tfrac12 (1-\lambda)^n \qquad \text{for all } n \geq 1 .
\end{equation}
Take $n = 1$: this gives $1 - \lambda' = \tfrac12 (1-\lambda)$. Square both sides:
\begin{equation}
(1-\lambda')^2 \;=\; \tfrac14 (1-\lambda)^2 .
\end{equation}
But taking $n = 2$ in the requirement gives $(1-\lambda')^2 = \tfrac12 (1-\lambda)^2$.
Since $\tfrac14 \neq \tfrac12$, no such $\lambda'$ exists.

The reason is structural: a mixture with a point mass at zero is not a geometric. The
tail of \eqref{eq:mixture} decays at rate $(1-\lambda)$, so matching the tail forces
$\lambda' = \lambda$, and then every non-zero length is off by exactly two.

\subsection{Why halving $\delta_s$ is not a fix either}

A subtler hope: leave the overhang wrong and halve $\delta_s$. Look
at~\eqref{eq:perpos} --- the leading term is a product $\Pr(\Ol \geq n)\,\delta_s$,
so halving $\delta_s$ does restore the per-position rate. This works, and it is worth
understanding why it is still not the same model.

It fails as soon as you ask about \emph{two} positions in the same read. Take
positions 1 and 2, both C in the reference, and ignore $\delta_d$:
\begin{align}
\text{correct model:} \qquad
  & \Pr(\Ol \geq 2)\, \delta_s^2 \;=\; \tfrac12 (1-\lambda)^2 \delta_s^2 , \\
\text{old model, } \delta_s \to \delta_s/2: \qquad
  & (1-\lambda)^2 \left(\tfrac{\delta_s}{2}\right)^{\!2}
    \;=\; \tfrac14 (1-\lambda)^2 \delta_s^2 .
\end{align}
Another factor of two, in the opposite direction. The marginals agree and the joint
does not. Intuitively: the correct model says half your molecules have no overhang at
all and the other half have an overhang whose Cs are deaminated with high probability
$\delta_s$; the rescaled model says every molecule has an overhang but each C is
deaminated only half as often. Those look identical one position at a time and quite
different at the level of a read --- which is precisely the level at which tools like
PMDtools and the read classification in ngsBriggs operate.

Measured on $600{,}000$ fragments with the actual binaries:
\begin{center}
\begin{tabular}{lccc}
\toprule
 & $\Pr(\text{pos } 1)$ & $\Pr(\text{pos } 2)$ & $\Pr(\text{pos }1 \text{ and } 2)$ \\
\midrule
correct model, $\delta_s = 0.68$            & 0.2244 & 0.1492 & 0.0947 \\
old model, $\delta_s$ halved to $0.34$      & 0.2207 & 0.1433 & 0.0464 \\
\bottomrule
\end{tabular}
\end{center}

\subsection{The fix}

\begin{verbatim}
    do{
        overhang5p = (randomProb(!seedSpecified)<0.5) ? overhang(generator) : 0;
        overhang3p = (randomProb(!seedSpecified)<0.5) ? overhang(generator) : 0;
    }while( (overhang5p+overhang3p) > (int(seq.size())-2) );
\end{verbatim}

\noindent
The rejection loop is the second half of the same paragraph in the paper: the two
overhangs together have to leave at least two bases of double-stranded DNA, otherwise
the molecule falls apart and is never sequenced. NGSNGS enforces the same condition.
The old code instead let the overhangs meet and then treated the whole fragment as
single-stranded --- negligible at $\lambda = 0.36$ on $60$\,bp fragments, but not for
short fragments or long mean overhangs.

\section{Correction 2: the nick position}

\subsection{What a nick does}

A nick is a break in the backbone of the \emph{opposite} strand. During library
preparation the strand-displacing polymerase extends from it, so everything downstream
of the nick in the final read is a copy of the opposite strand rather than the
original. Consequence: upstream of the nick a deaminated C reads as C$\to$T, and
downstream of it the opposite strand's deaminated Cs read as G$\to$A. This is what
produces the slow crossover from C$\to$T to G$\to$A along the interior of the
molecule, and \texttt{deamSim} models it correctly.

What was wrong is \emph{where} the nick goes.

\subsection{Uniform, not geometric}

Nicks arise at rate $\nu$ per base, so naively the first nick is geometric. Two things
break that:

\begin{enumerate}
\item A nick sitting in a single-stranded overhang is not observable. It does not
      split anything --- it just makes the overhang shorter. Only the double-stranded
      region can carry an observable nick.
\item If a fragment has a nick on one strand and a second nick further along on the
      other strand, the two replication forks meet during nick repair and the molecule
      splits in two. That fragment is lost and never sequenced.
\end{enumerate}

Point 2 removes the molecules with many nicks, and the surviving ones are those with
essentially one. Briggs et al.\ state the outcome directly: this ``causes the
distribution of first nicks in the sequenced fragments to be uniform rather than
geometric''.

Concretely, let $\ell$ and $r$ be the two overhang lengths and let
$D = L - \ell - r$ be the number of double-stranded bases. Let $M$ be the last
position still read from the original strand, so C$\to$T applies at $\delta_d$ for
$i \leq M$ and G$\to$A at $\delta_d$ for $i > M$. Then
\begin{equation}
\Pr(M = m) = \frac{\nu}{Z} \quad \text{for each of the } D-1 \text{ interior positions},
\qquad
\Pr(\text{no nick}) = \frac{1-\nu}{Z} ,
\end{equation}
with the normalizing constant
\begin{equation}
Z \;=\; (D-1)\,\nu \;+\; (1 - \nu) .
\end{equation}
Uniform over where the nick can be, with the leftover mass on ``no nick at all''.

\medskip
\noindent
\textbf{What \texttt{deamSim} did instead.} It walked the fragment from position 1,
flipping a $\nu$-coin at each base and stopping at the first success --- a geometric
--- and it did that walk across the overhangs too, where by point~1 above a nick
should not count.

The two disagree most visibly on how often there is no nick at all. At $L = 60$ and
$\nu = 0.024$:
\begin{equation}
\underbrace{(1-\nu)^{L} = 0.976^{60} = 0.233}_{\text{old}}
\qquad \text{versus} \qquad
\underbrace{\frac{1-\nu}{59\nu + 1 - \nu} = \frac{0.976}{2.392} = 0.408}_{\text{Briggs}} .
\end{equation}

\subsection{How much this matters}

Much less than correction 1, because everything here is multiplied by
$\delta_d \approx 0.01$ rather than $\delta_s \approx 0.68$. Its signature is the
position of the C$\to$T / G$\to$A crossover in the interior of the molecule --- the
geometric put it too far towards the $5'$ end. Simulated on $600{,}000$ $60$\,bp
fragments:

\begin{center}
\begin{tabular}{lcccc}
\toprule
& \multicolumn{2}{c}{corrected} & \multicolumn{2}{c}{old} \\
\cmidrule(lr){2-3}\cmidrule(lr){4-5}
position & C$\to$T & G$\to$A & C$\to$T & G$\to$A \\
\midrule
15 & 0.0088 & 0.0014 & 0.0069 & 0.0032 \\
25 & 0.0072 & 0.0027 & 0.0052 & 0.0045 \\
35 & 0.0061 & 0.0036 & 0.0040 & 0.0056 \\
45 & 0.0053 & 0.0048 & 0.0033 & 0.0075 \\
\bottomrule
\end{tabular}
\end{center}

\noindent
Note that the corrected version keeps C$\to$T above G$\to$A for much longer, which is
what the paper reports seeing in the data: ``the frequency of C to T substitutions
decreases steadily throughout the molecule toward $3'$ ends \ldots\ whereas G to A
substitutions do not seem to be elevated to the very $5'$ ends''. Since $\nu$ is
estimated largely from this crossover, this is the most likely reason mapDamage
recovered a wrong nick frequency from gargammel output.

\section{Does the fix work?}

Three independent things to check against.

\paragraph{Against the paper's own data.} Briggs et al.\ report $21\%$ C$\to$T on the
$5'$-most base of the Neandertal reads, and fit their four parameters to reproduce
that. Feeding those same fitted parameters back through the simulator has to return
roughly $21\%$.

\paragraph{Against an independent implementation.} NGSNGS
(\texttt{Biotin\_ds\_454Roche} in \texttt{Briggs.cpp}) implements the same model from
scratch, in a different group, with the $\tfrac12$ and the uniform nick.

\paragraph{Against the old behavior.} \texttt{-damagelegacy} has to reproduce the
previous output exactly, so that old simulations remain reproducible. (It does: the
output is byte-identical for a given seed.)

\begin{center}
\begin{tabular}{ccccc}
\toprule
position & corrected \texttt{-damage} & \texttt{-damagelegacy} & NGSNGS & Briggs 2007 observed \\
\midrule
1 & 0.2244 & 0.4382 & 0.2254 & 0.21 \\
2 & 0.1492 & 0.2805 & 0.1469 & \\
3 & 0.0966 & 0.1863 & 0.0979 & \\
4 & 0.0658 & 0.1199 & 0.0653 & \\
5 & 0.0456 & 0.0825 & 0.0460 & \\
6 & 0.0322 & 0.0536 & 0.0319 & \\
\bottomrule
\end{tabular}
\end{center}

\noindent
All simulated with $\nu = 0.024$, $\lambda = 0.36$, $\delta_d = 0.0097$,
$\delta_s = 0.68$ on $600{,}000$ random $60$\,bp fragments.

The corrected column agrees with the independent implementation to within Monte Carlo
noise, and lands close to the observed $21\%$; the legacy column is twice too large.
The remaining gap between $0.224$ and $0.21$ is the real data being slightly below the
fitted model, not a property of the simulator.

\section{Summary}

\begin{center}
\begin{tabular}{lll}
\toprule
 & old \texttt{deamSim} & Briggs et al.\ 2007 \\
\midrule
$\Pr(\Ol \geq n)$ & $(1-\lambda)^n$ & $\tfrac12 (1-\lambda)^n$ \\[2pt]
overhangs may meet & yes, whole fragment ss & no, rejected \\[2pt]
first nick & geometric, whole fragment & uniform, double-stranded part only \\[2pt]
$5'$-most C$\to$T at the MLEs & $0.438$ & $0.224$ \\
\bottomrule
\end{tabular}
\end{center}

\noindent
The first row is the whole story; everything else is a detail by comparison.

\section*{References}

\begin{itemize}
\item Briggs, A.W., Stenzel, U., Johnson, P.L.F., Green, R.E., Kelso, J., Pr\"ufer, K.,
      Meyer, M., Krause, J., Ronan, M.T., Lachmann, M., P\"a\"abo, S. (2007).
      Patterns of damage in genomic DNA sequences from a Neandertal.
      \emph{PNAS} 104(37):14616--14621. Model description on p.\ 14619 and in the SI Text.
\item Zhao, L., Henriksen, R.A., Rams\o{}e, A., Nielsen, R., Korneliussen, T.S. (2025).
      Revisiting the Briggs ancient DNA damage model: a fast maximum likelihood method
      to estimate post-mortem damage. \emph{Molecular Ecology Resources} 25:e14029.
\item NGSNGS, \texttt{Briggs.cpp}, function \texttt{Biotin\_ds\_454Roche}:
      \url{https://github.com/RAHenriksen/NGSNGS}
\end{itemize}

\end{document}
